\section{Planning with Diffusion}

In this section we describe a denoising diffusion model designed for trajectory data.
To plan with a diffusion model, we need to be able to condition its sampled trajectories on both evidence -- such as current state -- and desired outcomes -- either in the form of objectives to maximize or goals to reach.
We discuss methods for planning with diffusion models for each of these problem settings and describe how they are analogous to well-known sampling routines in deep generative modeling.

sampling strategies well known in deep generative modeling.
conditional sampling strategies used to conditionally sample from generativewell-known in the gen
We then discuss two options for planning with diffusion models.
We then turn to how to plan with the model, and elaborate on two possible options.
Both options mirror strategies well-known in generative modeling of images,
used for class-conditional sampling from an unconditional model and for image inpainting.
We do so by the use of two analogies to well-known procedure in the generative modeling of images: ''' '''.
We refer to the model and associated planning procedures as \method.

\subsection{A Generative Model of Trajectories}
We focus on trajectory-centric representations of data.
Instead of building a generative model for single-step prediction, we build one designed to produce an entire trajectory.
The data representation is a sequence of state action pairs: $\tau = [[s_1, a_2], [s_2, a_2]]$.

The architecture used to parameterize the reverse process $q_\theta$ is a U-Net structure with repeated blocks of residual temporal connections and group normalization (Figure~\ref{fig:architecture}.
The model weights depend on the dimensions of the states and actions, but not on the horizon of the plan; this is determined according to the input noise.

It is important to note what this model is not: it is not autoregressive in the planning horizon.
It iteratively refines the entire trajectory simultaneously.

%% MJ.1.20 There are a few options for the objective, depending on whether the model predicts the noise or the posterior process mean. I'm writing it up as the more conventional version -- predicting the noise -- but might swap it out later depending on final results. (It would be a local change that wouldn't affect anything outside of the next paragraph.)
Training


\subsection{Probabilistic Trajectory Optimization as Guided Sampling}
Training the diffusion model on large amounts of data acts as an unconditional prior for the types of behavior that the model knows about.
However, in many applications, we would like to steer the model toward particular regions.
To do so, we use auxiliary optimality signals.

%% Derivation

This amounts to a $Q$-value guided sampling analogous to classifier-guided sampling for image generation.
The source of the $Q$-values is independent of the method; and in particular any functional form could be used for gradients to steer the function.
We describe the training of a denoised $Q$-function in Algorithm~\ref{alg:q_training}.
We can use both Monte Carlo estimates from the data or train a $Q$-function using a reinforcement learning algorithm.


\subsection{Collocation as Inpainting}
An alternative control setting is not when we want to maximize a function, but when we want to sample any trajectory consistent with some constraints.
We can accomplish this by modifying the posterior distributions during training to reveal the conditioning information.
If the data representation is thought of as a single-channel image, this conditional sampling strategy is analogous to image inpainting.

Conditioning can in principle happen on any combination of the inputs, including non-contiguous regions of states and actions.
A particularly natural setting is that when a future state is revealed; this corresponds to sampling a trajectory consistent with the dynamics of the environment that begins at the current state and reaches a future state at a specific time.

There are two differences between the strategies discussed here: whether a function is to be maximized or evidence is to be conditioned on, and whether that evidence needs to be a part of the observation or action itself.

\section{Receding-Horizon Control with Diffusion Warm-Starting}

%%%%%%%%%%%%%%%%%%%%%%%%%%%%%%%%%%%%%%%%%%%%%%%%%%%%%%%%%%%%%%%%

\section{Properties of Diffusion Planners}

While \method can be thought of as a variant of a dynamics model, it has a different set of affordances than regular single-step models.
In this section, we describe some of the properties that are consequences of the design.

\paragraph{Non-greedy planning}
In contrast to standard single-step models, which can be incorporated directly into any classical planning algorithm,
the form of the planner is largely determined by the form of the model.
As a result, the model is not only an effective long-horizon predictor, but also an effective long-horizon planner.
We demonstrate the reason in a simplified 2D environment, in which an agent must plan around an obstacle to reach a goal location.
While shooting-style methods struggle because of the lack of reward shaping, our method behaves like a learned collocation technique.
We explore a more quantitative version of this problem setting in Section~\ref{sec:maze2d}.

\paragraph{Temporal compositionality}
A standard reason to use dynamics models is that they are the minimal component required to tackle an MDP problem.
Trajectory-level models lose out on the temporal compositionality that one can leverage due to the Markov property.
We show that \method also benefits from this type of compositionality.
In Figure~\ref{fig:properties}b, we train a diffusion model on trajectories that only exist on the right or left part of a state space.
Because the model uses local coherence to enforce global consistency, it can construct a trajectory that spans the entire image.

\pararagraph{Cost compositionality}
An important property of \method is that the cost function is not baked in by the training procedure, in contrast to model-free $Q$-functions and policies.
Because the model acts as a large unconditional prior over possible behaviors, it is possible to steer the model with comparatively lightweight cost functions at test-time.
We demonstrate this by adding hazards to the toy environment of the running example.
The model has remained unchanged, but its plans have updated to reflect a new test-time cost function.

\paragraph{Variable-length plans}
The minimal property of single-step models allows the planning horizon to be dynamically chosen at testing time; it need not be set in training time.
Because our model is fully convolutional in the horizon dimension, the same holds true here.
In Figure~\ref{fig:properties}d, we show how passing noise inputs of varying dimensionality to the model cause it to generate plans of varying horizon.


%%%%%%%%%%%%%%%%%%%%%%%%%%%%%%%%%%%%%%%%%%%%%%%%%%%%%%%%%%%%%%%%

\section{Experiments}

The aim of our experiments is to evaluate \method in data-driven control scenarios.
We begin in standard offline RL benchmark environments where conventionally model-free policies are the strongest.
In contrast, prior planning-based approaches have been limited to TT and MBOP.
This setting allows us to evaluate the guided sampling style of trajectory optimization.

We then turn our attention to settings requiring more test-time flexibility than standard RL benchmark environments, including 
a long horizon, sparse reward maze environment, for which we use the conditional inpainting version of planning, and a block stacking
ask for which test and train conditions are not identical.

\subsection{Offline reinforcement learning}
%% D4RL with Monte-Carlo Q-value guidance
%% D4RL with IQL Q-value guidance

\subsection{Goal-conditioned reinforcement learning}
%% Maze2D with inpainting

\subsection{Flexible behavior}
%% block stacking

%%%%%%%%%%%%%%%%%%%%%%%%%%%%%%%%%%%%%%%%%%%%%%%%%%%%%%%%%%%%%%%%

%% 1. method
%% 2. properties
%% 3. experiments

%%%% figures
%% 1. architecture
%% 2. teaser (sawyer)
%% 3. properties
%% 4. block stacking (kuka)

%%%% experiments
%% 1. goal-conditioned IQL for maze2d
%% 2. antmaze
%% 3. IQL Q-function guided locomotion

